% Part V -- Structured and Unstructured Data
% Chapter 14: Structured and unstructured data

\h{Structured and Unstructured Data}
\label{ch:14}

Every chapter so far has assumed the data already arrived as a table. Real
engineering projects rarely start there. A building performance study might
draw on a sensor log that is already a clean table, an asset record stored as
nested JSON, a stack of maintenance work orders written in prose, and a folder
of thermal photographs. These are not four difficulties of the same kind. They
differ in one specific way -- where the structure comes from -- and that
difference determines what you have to do before any analysis can begin.

\begin{NoteBox}
\textbf{Prerequisites.} You should be comfortable with dictionaries and lists
from \booklink{ch:03}{Chapter 3}, reading files and JSON from
\booklink{ch:10}{Chapter 10}, NumPy arrays from
\booklink{ch:11}{Chapter 11}, and pandas DataFrames from
\booklink{ch:12}{Chapter 12}. All data in this chapter is small and written
directly in the listings, so nothing needs to be downloaded.
\end{NoteBox}

\begin{tcolorbox}[title=Learning outcomes,colframe=orange,colback=white,arc=2mm]
After completing this chapter, you should be able to:
\begin{itemize}
    \item place a data source on the structured-to-unstructured spectrum;
    \item explain what "unstructured" actually means, and what it does not mean;
    \item state which schema a structured source must satisfy, and validate it;
    \item flatten a nested JSON record with \c{pd.json_normalize()};
    \item use \c{record_path} and \c{meta} to expand a repeating list;
    \item convert free text into a numeric feature table you have defined;
    \item describe an image as a NumPy array and reduce it to a feature row;
    \item explain why every path ends in one rectangular feature table; and
    \item state the judgment and information loss that feature extraction involves.
\end{itemize}
\end{tcolorbox}

\hh{Structure is a spectrum, not a binary}

The useful question about a data source is not "is it structured?" but "who
supplies the schema, and how much work is left for me?"

\bookfigure[0.98\textwidth]{Figures/Generated/ch14_data_spectrum.pdf}
{Structured files carry their own schema, semi-structured files describe themselves irregularly, and unstructured payloads carry no schema at all, so the work of imposing structure shifts steadily onto the analyst.}
{fig:ch14-spectrum}

\begin{center}
\begin{tabular}{lll}
\textbf{Class} & \textbf{Where the schema lives} & \textbf{Building example} \\
Structured & in the file, fixed and known & BMS sensor log, meter export \\
Semi-structured & in the file, but irregular & asset record JSON, IFC, GeoJSON \\
Unstructured & nowhere; you define it & work orders, photos, point clouds \\
\end{tabular}
\end{center}

\begin{ExplainBox}{"Unstructured" is a statement about the file, not the content}
A maintenance work order is not meaningless. It has grammar, vocabulary,
conventions, and a great deal of meaning that a technician reads instantly. What
it lacks is a schema that a program can rely on: no guaranteed fields, no
declared types, no fixed number of values per record. The same is true of a
photograph, which is highly structured as a grid of pixels but carries no
statement of what those pixels mean.

So "unstructured" means the structure needed for analysis is not supplied by the
storage format. Somebody has to impose it, and in a data project that somebody
is you. This matters because the choice is then yours to defend: two analysts
can extract different features from the same photographs and reach different
conclusions, and neither file format will tell you which was right.
\end{ExplainBox}

\hh{All three paths end at the same destination}

Whatever the source, analysis needs a rectangular table: one row per sample, one
column per feature. That is what pandas operates on, what the charts in
\booklink{ch:13}{Chapter 13} consume, and exactly the $X$ matrix that
\booklink{ch:19}{Chapter 19} and \booklink{ch:21}{Chapter 21} require.

\bookfigure[0.98\textwidth]{Figures/Generated/ch14_handling_paths.pdf}
{Structured sources need validation, semi-structured sources need flattening, and unstructured sources need feature extraction, but all three converge on one rectangular feature table.}
{fig:ch14-paths}

The three classes therefore do not need three different mindsets. They need
three different amounts of work to reach the same shape.

\hh{Structured data: state the schema and check it}

Structured data arrives with columns and types already decided. The risk is not
that you cannot read it; it is that you assume it is correct.

\begin{lstlisting}[
    language=python,
    style=block,
    caption={A structured sensor log},
    label={c14_structured},
    captionpos=b
]
import pandas as pd

readings = pd.DataFrame({
    "timestamp": pd.to_datetime([
        "2026-03-01 08:00", "2026-03-01 09:00", "2026-03-01 10:00",
        "2026-03-01 11:00", "2026-03-01 12:00",
    ]),
    "room_id": ["R101", "R101", "R101", "R101", "R101"],
    "temp_c": [20.4, 21.1, 24.8, 22.0, None],
    "co2_ppm": [612, 705, 1180, 890, 830],
})

print(readings)
print("rows, columns:", readings.shape)
print("missing values:", int(readings.isna().sum().sum()))
\end{lstlisting}

Expected output:

\begin{lstlisting}[language=text, style=block, caption={Clean shape, one gap}, label={c14_structured_output}, captionpos=b]
            timestamp room_id  temp_c  co2_ppm
0 2026-03-01 08:00:00    R101    20.4      612
1 2026-03-01 09:00:00    R101    21.1      705
2 2026-03-01 10:00:00    R101    24.8     1180
3 2026-03-01 11:00:00    R101    22.0      890
4 2026-03-01 12:00:00    R101     NaN      830
rows, columns: (5, 4)
missing values: 1
\end{lstlisting}

Being structured did not make the data complete: one temperature reading is
missing. Structure guarantees a shape, never a fact.

\hhh{Validate against a stated schema}

A \textbf{schema} is a written statement of which columns must exist and what
kind of value each holds. Checking it turns a silent wrong answer into a loud
error at the moment the data arrives.

\begin{lstlisting}[
    language=python,
    style=block,
    caption={Validate column presence and type kind},
    label={c14_validate},
    captionpos=b
]
import pandas as pd

EXPECTED = {
    "timestamp": "datetime",
    "room_id": "string",
    "temp_c": "number",
    "co2_ppm": "number",
}

CHECKS = {
    "datetime": pd.api.types.is_datetime64_any_dtype,
    "string": lambda column: (
        pd.api.types.is_string_dtype(column)
        or pd.api.types.is_object_dtype(column)
    ),
    "number": pd.api.types.is_numeric_dtype,
}


def validate_schema(frame, expected):
    """Return a list of schema problems; an empty list means the frame passed."""
    problems = []
    for name, kind in expected.items():
        if name not in frame.columns:
            problems.append(f"missing column: {name}")
        elif not CHECKS[kind](frame[name]):
            problems.append(f"{name} is {frame[name].dtype}, expected {kind}")
    return problems


print("as loaded:", validate_schema(readings, EXPECTED))

damaged = readings.astype({"co2_ppm": "str"})
print("after damage:", validate_schema(damaged, EXPECTED))
\end{lstlisting}

Expected output:

\begin{lstlisting}[language=text, style=block, caption={A schema check that catches a real defect}, label={c14_validate_output}, captionpos=b]
as loaded: []
after damage: ['co2_ppm is str, expected number']
\end{lstlisting}

\begin{SyntaxBox}{Check the kind of type, not its exact name}
It is tempting to compare \c{str(frame[name].dtype)} against a literal such as
\c{"int64"}. Resist it. The exact dtype string depends on the pandas and NumPy
versions, the platform's integer width, and whether missing values forced a
column to float. A check written that way fails on a colleague's machine for
reasons that have nothing to do with the data.

The \c{pd.api.types} predicates ask the question you actually mean -- is this
numeric, is this a datetime -- and stay correct across versions. Returning a
list of problems rather than raising immediately lets the caller report every
defect at once instead of fixing them one run at a time.
\end{SyntaxBox}

\begin{TryBox}{Make the check earn its place}
Delete the \c{co2_ppm} column entirely and rerun \c{validate_schema}. Then add
an unexpected extra column and rerun it again. Explain why the second case
reports no problem, and decide whether that is the right behavior for a study
whose input files come from an instrument you do not control.
\end{TryBox}

\hh{Semi-structured data: flatten what describes itself}

Semi-structured data carries its field names along with its values, so it can
change shape from record to record. JSON is the common case; IFC building
models, GeoJSON, and XML behave the same way. The difficulty is nesting: a value
may itself be an object or a list, and a DataFrame column cannot usefully hold
either.

\begin{lstlisting}[
    language=python,
    style=block,
    caption={A nested asset record},
    label={c14_nested},
    captionpos=b
]
asset = {
    "asset_id": "AHU-3",
    "location": {"building": "Link Hall", "floor": 2, "room": "M204"},
    "commissioned": "2019-06-14",
    "filters": [
        {"stage": 1, "type": "MERV8", "last_changed": "2026-01-10"},
        {"stage": 2, "type": "MERV13", "last_changed": "2025-11-02"},
    ],
}

print("top-level keys:", list(asset))
\end{lstlisting}

This one record mixes three shapes: plain values, a nested object
(\c{location}), and a list of objects (\c{filters}). Each needs different
handling.

\bookfigure[0.98\textwidth]{Figures/Generated/ch14_json_flatten.pdf}
{A nested JSON asset record becomes a rectangular table when json\_normalize expands the repeating filter list and carries the parent fields down to every row.}
{fig:ch14-flatten}

\hhh{Flatten nested objects with a separator}

\begin{lstlisting}[
    language=python,
    style=block,
    caption={Nested objects become dotted column names},
    label={c14_normalize},
    captionpos=b
]
import pandas as pd

flat = pd.json_normalize(asset, sep=".")

print("columns:", list(flat.columns))
print("floor value:", flat.loc[0, "location.floor"])
\end{lstlisting}

Expected output:

\begin{lstlisting}[language=text, style=block, caption={Objects flatten; lists do not}, label={c14_normalize_output}, captionpos=b]
columns: ['asset_id', 'commissioned', 'filters', 'location.building', 'location.floor', 'location.room']
floor value: 2
\end{lstlisting}

Read that output carefully. The nested \c{location} object became three ordinary
columns, but \c{filters} is still a single column holding a Python list. Nesting
of objects flattened; nesting of \emph{records} did not.

\hhh{Expand a repeating list into rows}

A list of records is not a wider table, it is a taller one. \c{record_path}
names the repeating list, and \c{meta} names the parent fields to copy onto
every resulting row.

\begin{lstlisting}[
    language=python,
    style=block,
    caption={One row per filter, with parent context preserved},
    label={c14_record_path},
    captionpos=b
]
filters = pd.json_normalize(
    asset,
    record_path="filters",
    meta=["asset_id", ["location", "building"]],
)

print(filters)
print("shape:", filters.shape)
\end{lstlisting}

Expected output:

\begin{lstlisting}[language=text, style=block, caption={Two filters become two rows}, label={c14_record_path_output}, captionpos=b]
   stage    type last_changed asset_id location.building
0      1   MERV8   2026-01-10    AHU-3         Link Hall
1      2  MERV13   2025-11-02    AHU-3         Link Hall
shape: (2, 5)
\end{lstlisting}

\begin{SyntaxBox}{Reading the two arguments that do the work}
\begin{itemize}
    \item \c{record_path="filters"} tells pandas that this list is the unit of
          observation, so the result has one row per filter rather than one row
          per asset.
    \item \c{meta=["asset_id", ["location", "building"]]} lists parent fields to
          repeat on every row. A plain string names a top-level key; a nested
          list of strings names a path down through nested objects.
    \item Without \c{meta}, the rows would carry no indication of which asset
          they belong to, which is the most common way this operation quietly
          destroys information.
\end{itemize}
Choosing \c{record_path} is choosing what one row means. That decision belongs
to your research question, not to the file.
\end{SyntaxBox}

\begin{ExplainBox}{Flattening always chooses a grain}
The original record describes one asset that has two filters. The flattened
table describes two filters that share one asset. Both are faithful; they answer
different questions. If you want to count assets, the second table will
double-count unless you aggregate first.

This is the central hazard of flattening, and it is not a pandas quirk -- it is
what happens whenever hierarchical information is forced into a rectangle. State
the grain of every table you produce: "one row per filter", "one row per
room-hour". \booklink{ch:12}{Chapter 12}'s grouping operations exist largely to move
between grains deliberately.
\end{ExplainBox}

\hh{Unstructured text: define the features yourself}

Free text carries no schema. To analyze it you must decide what about it
matters, and that decision is the analysis.

\begin{lstlisting}[
    language=python,
    style=block,
    caption={Turn work orders into a numeric feature table},
    label={c14_text_features},
    captionpos=b
]
import re

import pandas as pd

work_orders = [
    "Occupant reports room too warm all afternoon; VAV box not responding.",
    "Chilled water valve leaking near AHU-3, water on floor.",
    "Room too cold in the morning, occupants using space heaters.",
    "Filter change overdue; airflow reduced and room warm.",
]

VOCABULARY = ["warm", "cold", "leak", "filter", "valve", "airflow"]


def text_features(documents, vocabulary):
    """Return one row per document with a flag per vocabulary term."""
    rows = []
    for text in documents:
        tokens = re.findall(r"[a-z]+", text.lower())
        row = {
            f"has_{word}": int(any(token.startswith(word) for token in tokens))
            for word in vocabulary
        }
        row["word_count"] = len(tokens)
        rows.append(row)
    return pd.DataFrame(rows)


features = text_features(work_orders, VOCABULARY)
print(features)
print("shape:", features.shape)
\end{lstlisting}

Expected output:

\begin{lstlisting}[language=text, style=block, caption={Four documents become four rows of numbers}, label={c14_text_features_output}, captionpos=b]
   has_warm  has_cold  has_leak  has_filter  has_valve  has_airflow  word_count
0         1         0         0           0          0            0          11
1         0         0         1           0          1            0           9
2         0         1         0           0          0            0          10
3         1         0         0           1          0            1           8
shape: (4, 7)
\end{lstlisting}

\begin{SyntaxBox}{Each line makes a decision you could defend differently}
\begin{itemize}
    \item \c{text.lower()} treats "Warm" and "warm" as the same word. That is
          usually right, and it destroys the distinction permanently.
    \item \c{re.findall(r"[a-z]+", ...)} keeps runs of letters and discards
          punctuation and digits, so the asset identifier \c{AHU-3} survives
          only as \c{ahu}. Every reference to a numbered unit collapses onto the
          same token.
    \item \c{token.startswith(word)} makes "leaking" match "leak". It would also
          make "filtered" match "filter", and "coldest" match "cold".
    \item \c{VOCABULARY} is the whole model. A symptom nobody put on that list
          is invisible to every downstream analysis.
\end{itemize}
\end{SyntaxBox}

\begin{ExplainBox}{Feature extraction is lossy, and that is the point}
Four sentences became twenty-eight numbers. You cannot reconstruct the sentences
from the table, and you were never meant to: the table keeps what you decided
was relevant and discards the rest, which is exactly what makes it comparable
across thousands of records.

The danger is forgetting that a decision was made. A report that says
"31 percent of work orders concern temperature" is really saying "31 percent
contained one of the words I chose, matched by the rule I wrote." Publish the
vocabulary and the matching rule alongside the number, and the claim becomes
checkable. This is the same discipline \booklink{ch:21}{Chapter 21} applies to
model selection: the result is only meaningful together with the procedure that
produced it.
\end{ExplainBox}

\hh{Unstructured images: an array you summarize}

An image needs no conversion to become numeric -- it already is a NumPy array.
What it needs is reduction: from thousands of pixels to a handful of features.

\begin{lstlisting}[
    language=python,
    style=block,
    caption={Build a small synthetic thermal panel},
    label={c14_image},
    captionpos=b
]
import numpy as np

rng = np.random.default_rng(7)

panel = np.zeros((48, 64, 3), dtype=np.int16)
panel[:, :, 0], panel[:, :, 1], panel[:, :, 2] = 40, 60, 90

rows, columns = np.ogrid[:48, :64]
hot_spot = ((rows - 14) ** 2 + (columns - 46) ** 2) < 90
panel[hot_spot] = [235, 175, 95]

panel = np.clip(panel + rng.integers(-8, 9, panel.shape), 0, 255)
panel = panel.astype(np.uint8)

print("shape:", panel.shape, "dtype:", panel.dtype)
print("one pixel:", panel[14, 46])
\end{lstlisting}

Expected output:

\begin{lstlisting}[language=text, style=block, caption={Height, width, and three color channels}, label={c14_image_output}, captionpos=b]
shape: (48, 64, 3) dtype: uint8
one pixel: [235 172  98]
\end{lstlisting}

\begin{SyntaxBox}{Reading an image's shape}
\c{(48, 64, 3)} is 48 rows of pixels, 64 columns, and 3 color channels. Two
details trip up beginners repeatedly. The first number is the height, not the
width, because arrays are indexed row-first. And \c{uint8} holds whole numbers
from 0 to 255, so arithmetic on it wraps around silently -- which is why the
noise is added while the array is the wider \c{int16} type and only converted
back after \c{np.clip()} has bounded the result.
\end{SyntaxBox}

\begin{lstlisting}[
    language=python,
    style=block,
    caption={Reduce an image to one row of features},
    label={c14_descriptor},
    captionpos=b
]
HOT_THRESHOLD = 140.0

grayscale = panel.mean(axis=2)

descriptor = {
    "mean_intensity": round(float(grayscale.mean()), 2),
    "max_intensity": round(float(grayscale.max()), 2),
    "hot_fraction": round(float((grayscale > HOT_THRESHOLD).mean()), 4),
}

print(descriptor)
\end{lstlisting}

Expected output:

\begin{lstlisting}[language=text, style=block, caption={3,072 pixels reduced to three numbers}, label={c14_descriptor_output}, captionpos=b]
{'mean_intensity': 73.06, 'max_intensity': 175.0, 'hot_fraction': 0.0928}
\end{lstlisting}

\bookfigure[0.98\textwidth]{Figures/Generated/ch14_image_to_features.pdf}
{The synthetic panel is a 48 by 64 by 3 array; thresholding its grayscale version isolates the warm region, and summarizing that region reduces 3,072 pixels to three comparable numbers.}
{fig:ch14-image}

\c{grayscale.mean()} averages the three color channels into one intensity per
pixel; \c{(grayscale > HOT_THRESHOLD).mean()} then exploits the fact that the
mean of a boolean array is the fraction of \c{True} values. The hot fraction of
0.0928 says just under ten percent of the panel exceeds the threshold.

\begin{ExplainBox}{The threshold is a modeling choice with consequences}
\c{HOT_THRESHOLD = 140.0} is not measured; it is chosen. Raise it and the hot
fraction falls; lower it and ordinary variation starts counting as a defect. The
number 0.0928 is meaningless without it, exactly as an optimization result is
meaningless without its objective in \booklink{ch:18}{Chapter 18}.

Note also what these three features cannot express. They do not record where the
warm region sits, whether it is one blob or many, or what shape it has -- a
photograph with the same hot fraction spread thinly across the whole panel would
produce identical numbers. If position or shape matters to your question, the
descriptor must be extended to capture it.
\end{ExplainBox}

\hh{Bringing the paths together}

Each source contributed differently, and they combine into one row describing
one room.

\begin{lstlisting}[
    language=python,
    style=block,
    caption={One row per room, assembled from three kinds of source},
    label={c14_combined},
    captionpos=b
]
import pandas as pd

combined = pd.DataFrame([{
    "room_id": "R101",
    "mean_temp_c": round(float(readings["temp_c"].mean()), 2),
    "max_co2_ppm": int(readings["co2_ppm"].max()),
    "filter_stages": int(filters.shape[0]),
    "warm_reports": int(features["has_warm"].sum()),
    "hot_fraction": descriptor["hot_fraction"],
}])

print(combined.to_string(index=False))
print("shape:", combined.shape)
\end{lstlisting}

Expected output:

\begin{lstlisting}[language=text, style=block, caption={The destination every path was heading toward}, label={c14_combined_output}, captionpos=b]
room_id  mean_temp_c  max_co2_ppm  filter_stages  warm_reports  hot_fraction
   R101        22.07         1180              2             2        0.0928
shape: (1, 6)
\end{lstlisting}

\c{readings["temp_c"].mean()} skips the missing value automatically, so
\c{mean_temp_c} is the mean of four readings, not five. That is a reasonable
default and a silent one; with many gaps you would want to record how many
values each average was based on.

\hh{Choosing a storage format}

\begin{center}
\begin{tabular}{lll}
\textbf{Format} & \textbf{Best for} & \textbf{Main limitation} \\
CSV & flat tables, sharing, archiving & no types, no nesting \\
JSON & nested records, APIs & verbose, no schema enforcement, whole-file reads \\
Parquet & large analytic tables & needs a library; not human-readable \\
SQLite & queries and relationships & schema must be designed first \\
\end{tabular}
\end{center}

Reach for a database once relationships and repeated queries matter more than
portability -- when several tables reference each other, when files exceed
memory, or when many people read the data concurrently. Below that threshold, a
documented CSV or Parquet file is easier to share and to reproduce.

\hh{Common mistakes and useful diagnoses}

\begin{itemize}
    \item \textbf{Assuming structured means correct.} A schema guarantees shape,
          not accuracy or completeness. Check for missing values and impossible
          readings.
    \item \textbf{Comparing dtypes as literal strings.} Use the
          \c{pd.api.types} predicates so the check survives a version change.
    \item \textbf{Flattening without stating the grain.} Record what one row
          means, and aggregate before counting entities.
    \item \textbf{Dropping \c{meta} when using \c{record_path}.} The expanded
          rows then lose all parent identity.
    \item \textbf{Treating a keyword vocabulary as neutral.} It is the model;
          publish it with the results.
    \item \textbf{Reporting extracted features without the extraction rule.}
          "31 percent mention temperature" is only checkable alongside the
          vocabulary and matching rule.
    \item \textbf{Confusing image height and width.} The array shape is
          \c{(height, width, channels)}.
    \item \textbf{Doing arithmetic on \c{uint8} pixels.} Values wrap around
          silently; widen the type first.
\end{itemize}

\hh{Practice ladder}

\hhh{Level 0 -- Read}

For each of these, name the class and the work required: a CSV of hourly meter
readings; a GeoJSON of building footprints; a folder of PDF inspection reports;
an IFC model exported from BIM software.

\hhh{Level 1 -- Modify}

Add \c{"heater"} to \c{VOCABULARY} and rerun \c{text_features}. Report which
work orders change, and explain why \c{"space heaters"} in the third order was
invisible before.

\hhh{Level 2 -- Explain}

The flattened \c{filters} table has two rows for one physical asset. Describe a
question that this table answers correctly and a question it answers wrongly,
and state the operation that fixes the second case.

\hhh{Level 3 -- Build}

Write \c{summarize_documents(documents, vocabulary)} that returns a one-row
DataFrame containing the document count, the mean word count, and the share of
documents matching each vocabulary term. Validate that the shares lie between 0
and 1.

\hhh{Level 4 -- Challenge}

Extend the image descriptor so it also reports the row and column of the
brightest pixel and the number of separate hot regions. Then construct a second
panel with the same \c{hot_fraction} but a different arrangement, and show that
your extended descriptor distinguishes them while the original does not.

\hh{Practice solutions}

\textbf{Level 0:} The meter CSV is structured and needs validation. The GeoJSON
is semi-structured and needs flattening, with geometry handled as in
\booklink{ch:15}{Chapter 15}. The PDF reports are unstructured and need text
extraction followed by feature definition. The IFC model is semi-structured:
deeply nested and self-describing, so it needs a chosen grain and flattening.

\textbf{Level 2:} The table answers "how many filter stages are of type MERV13?"
correctly, because its grain is one row per filter. It answers "how many assets
are in Link Hall?" wrongly, returning two instead of one, because the asset is
repeated once per filter. Aggregating with
\c{filters.groupby("asset_id")} before counting restores the asset grain.

\textbf{Level 3}

\begin{lstlisting}[
    language=python,
    style=block,
    caption={One reusable document summary},
    label={c14_solution_summary},
    captionpos=b
]
import pandas as pd


def summarize_documents(documents, vocabulary):
    if len(documents) == 0:
        raise ValueError("documents must not be empty")

    features = text_features(documents, vocabulary)
    summary = {
        "document_count": len(documents),
        "mean_word_count": round(float(features["word_count"].mean()), 2),
    }
    for word in vocabulary:
        summary[f"share_{word}"] = round(float(features[f"has_{word}"].mean()), 4)

    shares = [value for name, value in summary.items() if name.startswith("share_")]
    if not all(0.0 <= share <= 1.0 for share in shares):
        raise ValueError("a share fell outside the range 0 to 1")

    return pd.DataFrame([summary])
\end{lstlisting}

The shares are means of columns holding only 0 and 1, so each is the fraction of
documents that matched. The validation is cheap and would catch a future change
that accidentally counted occurrences instead of flagging presence.

\hh{Chapter checkpoint}

\begin{enumerate}
    \item What distinguishes structured, semi-structured, and unstructured data?
    \item What does "unstructured" actually claim about a file?
    \item What does a schema guarantee, and what does it not guarantee?
    \item Why check the kind of a dtype rather than its exact name?
    \item What does \c{record_path} decide?
    \item What does \c{meta} preserve, and what breaks without it?
    \item What is the grain of a table, and why must it be stated?
    \item Why is feature extraction from text necessarily lossy?
    \item What does the shape \c{(48, 64, 3)} describe?
    \item What must be published alongside any extracted-feature result?
\end{enumerate}

\hh{Checkpoint answers}

\begin{enumerate}
    \item Where the schema lives: fixed in the file, self-described but
          irregular in the file, or nowhere at all.
    \item That the storage format supplies no schema for the payload, so the
          analyst must impose one. It says nothing about whether the content is
          meaningful.
    \item It guarantees which columns exist and what types they hold; it
          guarantees nothing about accuracy, completeness, or plausibility.
    \item Exact dtype names vary with library version and platform, so a literal
          comparison fails for reasons unrelated to the data.
    \item Which repeating list becomes the unit of observation, and therefore
          what one row of the result means.
    \item Parent fields copied onto every expanded row; without it the rows lose
          all identification of which parent record they came from.
    \item What one row represents. Without it, counts and averages silently
          answer a different question than the one asked.
    \item Because it keeps only the aspects you chose to measure and discards
          everything else, which is what makes records comparable.
    \item 48 pixel rows, 64 pixel columns, and 3 color channels: height first,
          then width.
    \item The extraction rule itself -- vocabulary, thresholds, and matching
          logic -- because the number is not checkable without it.
\end{enumerate}

\begin{tcolorbox}[title=Part V summary,colframe=orange,colback=white,arc=2mm]
Data sources differ in where their structure comes from, not in whether they
have any. Structured files carry a schema you should still validate;
semi-structured files describe themselves but must be flattened to a stated
grain; unstructured payloads carry no schema, so you define the features and own
that choice. All three paths converge on one rectangular feature table, which is
what every later chapter consumes. Because the last path is a judgment rather
than a conversion, publish the vocabulary, thresholds, and rules that produced
the numbers. \booklink{ch:15}{Chapter 15} takes up a data type that sits between
these classes: geometry, which is semi-structured in storage and demands its own
coordinate discipline.
\end{tcolorbox}
